\section{Введение}
\label{sec:Section0} \index{Section0}

В современном мире постоянно растёт число киберугроз и атак на информационные системы. По данным исследования глобального центра исследований и анализа угроз Kaspersky GReAT []
2025 года, количество атак на финансовые системы выросло на 43\% по сравнению с 2024 годом в России. Сетевая безопасность является одним из ключевых элементов защиты информационных систем от киберугроз.

В последние 20 лет активно развивается рынок средств защиты информационных систем. Компании развивают не только решения для конечных пользователей (end-point security), но и решения для сетевого уровня (network security). Средства защиты применяются не только на привычных сетевых устройствах: роутерах или свитчах, но и применяются специализированные отдельные сетевые устройства: межсетевые экраны (firewalls). Межсетевые экраны обеспечивают защиту сети в соответствии с политикой безопасности. Под политикой безопасности сетевого устройства понимается формализованное описание того, какой трафик и при каких условиях является допустимым с точки зрения безопасности организации. Политика может охватывать не только фильтрацию трафика, но и смежные механизмы: трансляцию сетевых адресов (NAT), SSL-инспекцию зашифрованных соединений и другие. Центральным элементом политики является упорядоченный набор правил фильтрации. Каждое правило состоит из условия — набора квалификаторов пакета или сессии, которым он должен соответствовать, — и действия (например, разрешить, запретить).
Реализация поиска подходящего правила может быть реализована по разному, но поведение совпадает с моделью <<first-match>>: применяется первое подходящее правило. Если ни одно правило не сработало, применяется действие по умолчанию — как правило, запрет.

В текущий момент можно выделить несколько видов межсетевых экранов, но границы между классами устройств на рынке невсегда однозначны.

\textbf{Классический межсетевой экран} реализует политику доступа c помощью правил фильтрации сетевого трафика, имеющие квалификаторы сетевого и транспортного уровней модели OSI []: диапазоны IP адресов, номеров протоколов и набора портов (для UDP и TCP). Такие решения применяют для разграничение сегментов сети и контроля доступом. Аналогичный механизм фильтрации трафика реализован в виде списков контроля доступа (Access Control List, ACL) на маршрутизаторах и коммутаторах, где ACL определяют, какой трафик разрешён или запрещён на конкретном интерфейсе. Однако классический межсетевой экран имеет значительный недостаток: его правило нельзя непосредственно связать с конкретным прикладным сервисом или пользователю (субъекту) в полной мере.

Под \textbf{универсальным шлюзом безопасности} (Unified Threat Management) традиционно понимают узел объединяющий на одном аппаратном комплексе или в рамках одного программного решения большое количество средств защиты: межсетевая фильтрация, системы обнаружения и предотвращения вторжений (IDPS), антивирусная и URL фильтрация, организация защищённых каналов (VPN), анализ электронной почты. Администрирование ведётся через единый интерфейс.
Для малых и средних организаций такой подход может быть более экономически выгодным и удобным, чем поддержка набора функционально близких, но технически раздельных систем.

\textbf{Межсетевой экран нового поколения} (Next Generation Firewall, NGFW) является приемником UTM и расширяет возможности классического межсетевого экрана: идентификация приложений, сервисов, пользователей. Данный тип экрана включает в себя все современные механизмы безопасности: глубокий анализ пакетов (Deep Packet Inspection, DPI) с контролем приложений, систему обнаружения и предотвращения вторжений (IDPS), SSL-инспецию и гибкое управление политиками безопасности.

\textbf{Межсетевой экран веб-приложений} (Web Application Firewall, WAF) специализированное устройство для защиты сервисов, работающих по протоколам HTTP и HTTPS: проверке подвергаются запросы и ответы с целью блокирования атак на типовые веб-уязвимости. Данное устройство не заменяет сетевой межсетевой экран, а дополняет его, так как разграничение доступа между подсетями и реализация единой периметровой политики безопасности остаются задачами классического МЭ или NGFW.

В России активно развивается рынок межсетевых экранов нового поколения, появляются решения от разных компаний: Kaspersky, Positive Technologies, Ideco и др. Появление новых отечественных NGFW во многом связано с уходом с рынка ряда зарубежных вендоров (Check Point, Fortinet, Palo Alto Networks) и дальнейшим курсом на импортозамещение, вследствие чего организации вынуждены переходить на российские решения. При этом импортозамещение затрагивает не только сегмент NGFW и других средств защиты, но и базовое сетевое оборудование, такое как маршрутизаторы и коммутаторы. Вместе с заменой устройств приходится переносить и их конфигурации, что порождает ряд трудностей.

Во-первых, модели правил у разных вендоров отличаются, и по составу квалифакторов в правилах, и способу их задания: одни оперируют набором сетевых квалификаторов (адреса, протокол, порты), другие поддерживают именованные объекты, группы адресов или набоы портов. Прямой перенос правил <<один-к-одному>> часто оказывается невозможным или приводит к семантически неэквивалентной политике: правила могут быть избыточными, перекрываться или, напротив, оставлять непокрытое множество трафика. Во-вторых, в процессе миграции легко допустить ошибки конфигурации, которые не проявляются сразу, но впоследствии приводят либо к нарушению работы необходимых сервисов и приложений, либо к появлению уязвимостей в инфраструктуре организации.

Отдельной проблемой является эксплуатация уже существующих решений. Многие организации продолжают использовать классические межсетевые экраны и устройства с ACL, введенные в эксплуатацию несколько лет назад или даже десятилетия назад. За время эксплуатации в таких политиках накапливается значительное число правил. Некоторые из них были добавлены <<временно>> и были забыты, другие заводились по заявке пользователей, которые уже не работают в организации. Также могут накапливаться и ошибки в политике безопасности: дублирующиеся правила, перекрывающиеся правила, затенения правил и т.д. Объём политик безопасности в крупных корпоративных сетях нередко исчисляется десятками тысяч правил фильтрации, что делает ручной аудит политик практически невозможным, а попытки внести изменения — очень рискованными.

Таким образом, задача автоматизированного анализа правил фильтрации — поиска неиспользуемых правил, конфликтов, избыточности и несогласованностей в политиках безопасности — остаётся актуальной не только при сопровождении эксплуатируемой инфраструктуры, но и при миграции на новые сетевые устройства. Разработка системы, способной формально верифицировать правила фильтрации, позволяет гарантировать корректность и полноту фильтрации трафика. Это позволит организациям повысить качество и надежность сетевой защиты, минизировать риски, возникающие при изменении политики безопасности, и снизить трудозатраты инженеров, эксплуатирующих сетевую инфраструктуру.

\newpage